\chapter{Considerações Finais}
\label{chap:conclusao}

\orientacao{Meia a uma página de texto, mais os quadros. Quatro funções: declarar o cumprimento retomando a descrição do Anexo I; sintetizar o que foi alcançado em cada frente, uma frase por frente; declarar o nível de maturidade no início e no fim, dizendo qual dimensão mudou e apontando a evidência; indicar a base deixada para a entrega seguinte.

Não traga informação nova, não repita a discussão e não termine com frase de efeito sobre o futuro da tecnologia.}

\section{Cumprimento da Macroentrega}
\label{sec:cumprimento}

A Macroentrega~\numeromacroentrega, referente a ``\ph{descrição da macroentrega, igual à do Anexo I}'', foi \ph{concluída / concluída em parte, com a parte faltante identificada}. A Tabela~\ref{tab:cumprimento} relaciona cada cláusula da descrição contratual ao ponto do relatório que a responde e à evidência que a sustenta.

\orientacao{Quebre a descrição do Anexo I em cláusulas verificáveis, uma por linha. Cláusula sem resposta é lacuna de entrega e não se resolve com redação: leve à coordenação antes de escrever. Resultado sem cláusula correspondente é resultado adicional: descreva-o depois de demonstrar o escopo contratado e identifique-o como adicional, sem que ocupe o lugar de item contratado.}

\begin{table}[htbp]
\centering
\caption{Cumprimento da descrição da macroentrega no Anexo I.}
\label{tab:cumprimento}
\estilotabela
\begin{tabularx}{\linewidth}{|P{4.0cm}|P{2.4cm}|X|C{1.9cm}|}
\hline
\rowcolor{reqcolor}\cab{Cláusula do Anexo I} & \cab{Resposta} & \cab{Evidência} & \cab{Situação} \\ \hline
\ph{cláusula 1} & Seção~\ref{sec:frente1_resultados} & Fig.~\ref{fig:exemplo} & \done \\ \hline
\rowcolor{reqlight}\ph{cláusula 2} & \ph{Seção} & \ph{Fig. ou Tab.} & \parc \\ \hline
\end{tabularx}
\end{table}

\section{Síntese por Frente}

\begin{itemize}
    \item \ph{Frente 1: uma frase com o artefato ou achado e o número que o caracteriza.}
    \item \ph{Frente 2: uma frase com o artefato ou achado e o número que o caracteriza.}
\end{itemize}

\section{Maturidade Tecnológica}
\label{sec:trl}

\orientacao{Escreva na forma ``o projeto atinge o nível X, comprovado por Y'', e não ``o projeto atingiu o nível X''. O valor declarado precisa ficar dentro da faixa prevista no Anexo I e ser o mesmo lançado no SRInfo.

A evidência cobre as três dimensões da norma ABNT NBR ISO 16290: escala do objeto, fidelidade do objeto e ambiente em que a função tecnológica é desenvolvida. Defina o ambiente antes de usar o rótulo: uma frase com o envelope de condições que o Anexo I exige (iluminação, terreno, velocidade, volume de dados, perfil de uso) e outra com o subconjunto em que o ensaio desta macroentrega ocorreu. A distância entre as duas é o que o nível declarado precisa suportar.

Declare o nível no início e no fim e diga qual dimensão mudou. Se nenhuma mudou, provavelmente não houve avanço de nível, e declarar TRL mantido é resultado legítimo. Em entrega de integração, vale o nível do componente menos maduro do caminho crítico, e convém nomeá-lo. Ensaio em bancada não sustenta ``ambiente relevante''.}

O projeto atinge o nível TRL~\ph{X}, comprovado por \ph{evidência}. Os níveis de maturidade são caracterizados por escala do objeto, fidelidade do objeto e ambiente, conforme a norma ABNT NBR ISO 16290~\cite{abnt16290}. O ambiente exigido pelo Anexo I é \ph{envelope de condições}; o ensaio desta macroentrega ocorreu em \ph{subconjunto dessas condições}.

\begin{table}[htbp]
\centering
\caption{Nível de maturidade tecnológica no início e no final da macroentrega.}
\label{tab:trl}
\estilotabela
\begin{tabularx}{\linewidth}{|P{2.6cm}|X|X|P{2.6cm}|}
\hline
\rowcolor{reqcolor}\cab{Dimensão} & \cab{Início (TRL~\ph{X})} & \cab{Final (TRL~\ph{Y})} & \cab{Evidência} \\ \hline
Escala do objeto & \ph{escala} & \ph{escala} & \ph{Fig. ou Tab.} \\ \hline
\rowcolor{reqlight}Fidelidade do objeto & \ph{fidelidade} & \ph{fidelidade} & \ph{Fig. ou Tab.} \\ \hline
Ambiente & \ph{ambiente} & \ph{ambiente} & \ph{Fig. ou Tab.} \\ \hline
\end{tabularx}
\end{table}

\section{Próximos Passos}

\orientacao{Direção técnica em duas ou três frentes, cada uma com a dependência que a condiciona. Sem cronograma, sem responsável e sem métrica: tarefa, dono e prazo pertencem à ata e à proposta da entrega seguinte, e escrevê-los aqui cria compromisso não contratado. Prefira ``o passo seguinte é ampliar a coleta'' a ``o sistema passará a generalizar para novas situações''.}

\ph{Base estabelecida por esta entrega e direção técnica da entrega seguinte, em verbo modesto.}
